
The user can bind his/her own C++ classes to graphs. These classes are
called labels. You can bind your labels to graphs, nodes, and
edges using graph-labels, node-labels or edge-labels. If you bind a
node-label or an edge-label to a graph, a prototype of the label is
bound to every node or edge in the graph and will be bound to
every node or edge that will be added to the graph in the future.
There are 4 basic label classes included in the library. {\bf
  FFGraphLabel}, {\bf FFNodeLabel} and {\bf FFEdgeLabel} which
are all subclasses of {\bf FFLabelClass}. All graph-labels are
subclasses of {\bf FFGraphLabel}, {\bf FFNodeLabel} for nodes and
{\bf FFEdgeLabel} for edges. When a graph is loaded or saved the
content of the labels may be loaded or saved as well. Whether a
label is saved/loaded depends on the implementation of the label.
It is even possible that only the labels of certain nodes or
edges are saved/loaded. When a graph is copied the contents of
the labels may be copied as well. Whether a label is copied
depends on the implementation of the label. Either all or no
labels of a certain type are copied.  Labels get informed of any
structural changes to the graph by events that are sent to the
labels. A label assigned to a graph must have a unique type. Only
one label class of each type can be assigned to a certain graph.
The type is determined by the result of the method {\bf const
  String\& FFLabelClass::type() const} which each label class
provides. The type may not correspond to the C++ type in any
way. The type of a label usually consists of an arbitrary string
preceded by the type of its superclass. The two strings are
separated by ":".

\subsection{Using labels}  
\label{using_labels}

Labels are subclasses of the classes {\bf FFGraphClass}, {\bf
FFNodeClass} or {\bf FFEdgeClass} which are subclasses of the class
{\bf FFLabelClass}. The library contains methods to add, access and
delete labels from graphs. 

To add labels to a graph use the methods

\begin{itemize}

\item{\bf  FFGraphClass::addGraphLabel(FFGraphLabel *label)}

\item{\bf FFGraphClass::addNodeLabel(FFNodeLabel *label)}

\item{\bf FFGraphClass::addEdgeLabel(FFEdgeLabel *label)}

\end{itemize}

{\bf label} must be created using the {\bf new}
operator\footnote{Unless the manual of a certain label says
  anything else}. After calling one of these methods {\bf label}
belongs to the graph library and will usually be freed by it
using the {\bf delete} operator. Copies of node-labels and
edge-labels will automatically be assigned to every node and edge
in the graph and to all nodes and edges that will be added to the
graph in the future.

The return value of these methods is the index of the added
label. This index is local to a graph and may differ from graph
to graph. You can access a label through its type or its index.
Access through index is recommended as it is much
faster\footnote{index: $O(1)$, type: $O(|labels|)$}. To access a label by
its index or type you can use one of the following methods:

\begin{itemize}

\item{\bf FFGraphLabel *FFGraphClass::getGraphLabel(int i)}
  
  Returns the graph label with index {\bf i}, or {\bf NULL} if
  the graph has no label with index {\bf i}.
  
\item{\bf FFGraphLabel *FFGraphClass::getGraphLabel(const
    String\& type, AVLSet\verb|<|FFGraphLabel *\verb|>|
    *set=NULL)}
  
  If {\bf set} {\bf == NULL} the label with type {\bf type} is
  returned, or {\bf NULL} if the graph has no graph label of type
  {\bf type}.
  
  If {\bf set} {\bf != NULL} all labels that contain {\bf type
    }in their type are added to {\bf set}. The first label found
  is returned. If no graph label-type contains {\bf type}, {\bf
    NULL} is returned.
  
\item{\bf const MPlex\verb|<|FFGraphLabel *\verb|>|\&
    FFGraphClass::getGraphLabels() const}
  
  Returns a MPlex containing all graph labels.
  
\item{\bf const MPlex\verb|<|FFNodeLabel *\verb|>|\&
    FFGraphClass::getNodeLabels() const}
  
  Returns a MPlex containing the prototypes of all node labels.
  
\item{\bf const MPlex\verb|<|FFEdgeLabel *\verb|>|\&
    FFGraphClass::getEdgeLabels() const}

  Returns a MPlex containing the prototypes of all edge labels.

\item{\bf FFNodeLabel *FFNodeClass::getLabel(int i)}
  
  Returns the node label with index {\bf i} or {\bf NULL} if the
  node has no label with index {\bf i}.

\item{\bf const MPlex\verb|<|FFNodeLabel *\verb|>|
*FFNodeClass::getLabels() const}

Returns a MPlex containing all labels bound to the node.

\item{\bf FFEdgeLabel *FFEdgeClass::getLabel(int i)}
  
  Returns the edge label with index {\bf i} or {\bf NULL} if the
  edge has edge label with index {\bf i}.

\item{\bf const MPlex\verb|<|FFEdgeLabel *\verb|>|
*FFEdgeClass::getLabels() const}

Returns a MPlex containing all labels bound to the node.

\end{itemize}

If you did not store the index you got by adding a label, you can
retrieve its index by using one of the following methods:

\begin{itemize}
  
\item{\bf int FFGraphClass::getGraphLabelIndex(const String\&
    type, AVLSet\verb|<|int\verb|>|*set=NULL) const}

If {\bf set == NULL} the index of the graph label with type {\bf
  type} is returned or {\bf -1}, if the graph has no graph label
of type {\bf type}.

If {\bf set != NULL} the indexes of all graph labels that contain
{\bf type }in their type are added to {\bf set}. The first index
found is returned. If no graph label type contains {\bf type},
{\bf -1} is returned.

\item{\bf int FFGraphClass::getNodeLabelIndex(const String\& type,
AVLSet\verb|<|int\verb|>|*set=NULL) const}

If {\bf set == NULL} the index of the node label with type {\bf
  type} is returned or {\bf -1}, if the graph has no node label
of type {\bf type}.

If {\bf set != NULL} the indexes of all node labels that contain
{\bf type }in their type are added to {\bf set}. The first index
found is returned. If no node label type contains {\bf type},
{\bf -1} is returned.

\item{\bf int FFGraphClass::getEdgeLabelIndex(const String\&
    type, AVLSet\verb|<|int\verb|>|*set=NULL) const}

If {\bf set == NULL} the index of the edge label with type {\bf
  type} is returned or {\bf -1}, if the graph has no edge label
of type {\bf type}.

If {\bf set != NULL} the indexes of all edge labels that contain
{\bf type }in their type are added to {\bf set}. The first index
found is returned. If no edge label type contains {\bf type},
{\bf -1} is returned.

\item{\bf int FFNodeClass::getLabelIndex(const String\& type,
    AVLSet\verb|<|int\verb|>|*set=NULL) const}

If {\bf set == NULL} the index of the node label with type {\bf
  type} is returned or {\bf -1}, if the node has no node label of
type {\bf type}.

If {\bf set != NULL} the indexes of all node labels that contain
{\bf type }in their type are added to {\bf set}. The first index
found is returned. If no node label type contains {\bf type},
{\bf -1} is returned.

\item{\bf int FFEdgeClass::getLabelIndex(const String\& type,
    AVLSet\verb|<|int\verb|>|*set=NULL) const}

If {\bf set == NULL} the index of the edge label with type {\bf
  type} is returned or {\bf -1}, if the edge has no edge label of
type {\bf type}.

If {\bf set != NULL} the indexes of all edge labels that contain
{\bf type }in their type are added to {\bf set}. The first index
found is returned. If no edge label type contains {\bf type},
{\bf -1} is returned.

\end{itemize}

If you want to change a label, you can get the label using one of
the above methods and change the label by calling its methods.

If you want to assign the contents of one label to another, you
can either

\begin{itemize}
  
\item Get the label and call the method {\bf void
    FFLabelClass::assign(const FFLabelClass*)} which every label
  provides. Of course you cannot be sure that the implementation
  of this method really does anything at all.


\end{itemize}

or

\begin{itemize}

\item
Use one of the following methods:

\begin{itemize}
  
\item{\bf void FFGraphClass::assignGraphLabel(const FFGraphLabel
    *label, int i=-1)}

calls the method {\bf void FFLabelClass::assign(const
  FFLabelClass*)} of the graph label with index {\bf i} and
argument {\bf label}. If {\bf i} == -1 the graph checks whether it is
labeled with a label of the same type as {\bf label} and calls
{\bf assign} if it finds one.

\item{\bf void FFGraphClass::assignGraphLabel(const
MPlex\verb|<|FFGraphLabel *\verb|>|\& set)}
        
For each label in {\bf set} the graph checks if it is labeled
with a label of the same type. If it finds one, the method {\bf void
  FFLabelClass::assign(const FFLabelClass*)} is called
with the label from {\bf set} as argument.

\item{\bf void FFGraphClass::setGraphLabel(FFGraphLabel *label,
    int index=-1, BOOL doFree=YES)}
  
  Replaces the graph label of index {\bf index} with {\bf label}.  If
  {\bf doFree == YES} the old label is deleted using the {\bf delete}
  operator.  If {\bf index == -1} the graphs checks whether it is
  labeled with a label of the same type as {\bf label} and does the
  replacement if it finds one. After calling this method {\bf label}
  belongs to the graph and will be freed by the graph using the {\bf
    delete} operator when it is no longer needed.

\item{\bf void FFGraphClass::setGraphLabel(const
    MPlex\verb|<|FFGraphLabel *\verb|>|\& set, BOOL doFree=YES)}

For each label in {\bf set} the method replaces the graph label
of the same type with the label from {\bf set}. If {\bf doFree ==
  YES} the old label is deleted using the {\bf delete} operator.
After calling this method all labels of {\bf set} belong to the
graph and will eventually be freed by the graph using the {\bf
  delete} operator.

\item Use one of the corresponding methods for {\bf FFNodeClass} and
  {\bf FFEdgeClass}.

\end{itemize}

\end{itemize}

If you no longer need a label, you can remove it from the graph
using one of the following methods:

\begin{itemize}
  
\item{\bf void FFGraphClass::delGraphLabel(int index, BOOL
    doFree=YES) \\
    void FFGraphClass::delGraphLabel(const String\& type,
    BOOL doFree=YES)}
  
  Deletes the graph label with index {\bf index} or type {\bf
    type} from the graph. If {\bf doFree == YES} the label is
  freed using the {\bf delete} operator.
  
\item{\bf void FFGraphClass::delEdgeLabel(int index, BOOL
    doFree=YES) \\
    void FFGraphClass::delEdgeLabel(const String\& type,
    BOOL doFree=YES)}
  
  Deletes the edge labels with index {\bf index} or type {\bf
    type} from all edges of the graph. If {\bf doFree == YES} the
  labels are freed using the {\bf delete} operator.
  
\item{\bf void FFGraphClass::delNodeLabel(int index, BOOL
    doFree=YES) \\
    void FFGraphClass::delNodeLabel(const String\& type,
    BOOL doFree=YES)}
  
  Deletes the node labels with index {\bf index} or type {\bf
    type} from all nodes of the graph. If {\bf doFree == YES} the
  labels are freed using the {\bf delete} operator.

\end{itemize}

Graph-labels can contain an inspector to display their
content. Inspectors are subclasses of {\bf FFInspector}. The kind of
information displayed by the inspector depends on the implementation
of the inspector. E.g. the inspector contained in {\bf FFTclWindow}
displays the result of the {\bf String FFLabelClass::unparse()} method
for graph, node or edge labels.

To determine whether a label has its own inspector, you can use the
method:

\begin{itemize}

\item{\bf BOOL FFGraphLabel::hasOwnInspector()}

\end{itemize}

To display an inspector you can use the method:

\begin{itemize}

\item{\bf void FFGraphLabel::showInspector()}

\end{itemize}

To update the display of the inspector you can use the method:

\begin{itemize}

\item{\bf void FFGraphLabel::updateInspector()}

\end{itemize}

\subsection{Writing labels}
\label{writing_labels}

If you want to use your own label, you have to write a subclass
of {\bf FFGraphLabel}, {\bf FFNodeLabel} or {\bf FFEdgeLabel}
with any additional data and methods you want. You then can add
it to the graph using one of the {\bf add...Label(...)} methods.
Labels get informed when the structure of the graph is changed.
To overwrite the method {\bf void event(labelEvent,
  labelEventStruct *)} is the easiest way to respond to these
changes to the graph. To enable the library to handle your label
you must implement some additional methods. The methods you must
implement are:

\begin{itemize}
  
\item Write a constructor for your class. In the constructor you
  have to call the method {\bf void FFLabelClass::\_setType(const
    String\& s)} with an arbitrary String {\bf s}. {\bf s} must
  not contain ":" and must be unique. Using the C++ type of the
  label class often is a good choice.

\item Depending on what kind of label you are writing, one of 

\begin{itemize}

\item{\bf FFGraphLabel *YourLabelClass::proto()}

\item{\bf FFNodeLabel *YourLabelClass::proto()}

\item{\bf FFEdgeLabel *YourLabelClass::proto()}

\end{itemize}

This method must return a new instance of your label class. The
new instance should be allocated using the {\bf new} operator. This
method is needed to assign labels to nodes, edges and copies of
graphs.

If you e.g. write a graph label usually writing something like

\bigskip

\begin{verbatim}
FFGraphLabel *YourLabelClass::proto()
{
  return new YourLabelClass;
}
\end{verbatim}

\bigskip

is sufficient. 

\end{itemize}

There are also some methods you might want to overwrite, but you
do not have to:

\begin{itemize}

\item{\bf BOOL FFLabelClass::isSaveable() const}
  
  If this method returns {\bf YES} the result of the method {\bf
    String FFLabelClass::unparse()} is stored in a file when a
  graph is saved. When a graph is loaded the method {\bf BOOL
    FFLabelClass::parse()} is called with the string saved
  before.
  
  If this method returns {\bf NO} the label is not stored in the
  file.  The results of the method may differ between two nodes
  or edges in the graph for the same label.

This method defaults to return {\bf NO}.

\item{\bf BOOL shouldCopy() const}
  
  If a graph gets copied and this method returns {\bf YES} the
  new graph is labeled with the label as well. The {\bf assign}
  method is called for the label in the new graph with this label
  as argument. For node and edge labels {\bf assign} is called
  for labels in the corresponding node or edge in the new graph.

This method defaults to return {\bf YES}.

\item{\bf BOOL isInspectable() const}
  
  Inspector classes are used to display the contents of
  a label on screen. Usually the result of calling the method
  {\bf String FFLabelClass::unparse()} is displayed. If this
  method returns {\bf YES} it can be displayed by an inspector.

This method defaults to return {\bf NO}.

\item{\bf BOOL compare(const FFLabelClass *label) const}

This method should return if {\bf label} is equal to this label. 

\item{\bf void assign(const FFLabelClass *label)}

  This method should assign the content of {\bf label} to {\bf this}.

\item{\bf BOOL parse(const String\& str)}
  
  This method should be able to initialize itself according to
  {\bf str}. {\bf str} is usually the result of a previous call
  to {\bf String FFLabelClass::unparse()}. This method is e.g.
  called when a graph is loaded from a file or by an inspector
  to display the content of the label.  The method should return
  whether it could successfully parse {\bf str}.

\item{\bf String unparse()}
  
  This method should return the content of the label as a {\bf
    String}. An instance of the label class should be able to
  initialize itself if the result of invoking {\bf unparse()} is
  passed to {\bf BOOL parse(const String\& str)}. This method is
  e.g. called when a graph is saved to a file or by an
  inspector.

\item{\bf void event(labelEvent kind, labelEventStruct *data)}
  
  Whenever there is a structural change to the graph, the {\bf
    event} method is invoked. {\bf kind} specifies the kind of
  event that happened and the union {\bf data} contains some
  additional information that varies between different kinds of
  events. {\bf event} is not always invoked in all labels, but
  there are some events that are e.g. only invoked in graph
  labels etc.

{\bf labelEventStruct} is defined as follows:

\bigskip

\begin{verbatim}
union {
  FFNode node;
  FFEdge edge;
  int index;
} labelEventStruct;
\end{verbatim}

\bigskip

The following kinds of events may occur:

\begin{itemize}

\item{\bf ADD\_NODE}
  
  Is invoked after a new node has been added to the graph. {\bf
    data\verb|->|node} contains the new node.
  
  Is delivered to all graph labels and to all node labels,
  including the new node.

\item{\bf ADD\_EDGE}
  
  Is invoked after a new edge has been added to the graph. {\bf
    data\verb|->|edge} contains the new edge.
  
  Is delivered to all graph labels and to all edge labels,
  including the new edge.

\item{\bf DEL\_NODE}
  
  Is invoked before a node is removed from the graph. {\bf
    data\verb|->|node} contains the node that will be removed.
  
  Is delivered to all graph labels and to all node labels,
  including the node that will be deleted.

\item{\bf DEL\_EDGE}
  
  Is invoked before an edge is removed from the graph. {\bf
    data\verb|->|edge} contains the edge that will be removed.
  
  Is delivered to all graph labels and to all edge labels,
  including the edge that will be deleted.

\item{\bf DEL\_GRAPH}
  
  Is invoked before a graph is deleted.
  
  Is delivered to all graph labels.

\item{\bf REV\_EDGE}
  
  Is invoked after the direction of an edge has changed.
  
  Is delivered to all labels of this edge. {\bf
    data\verb|->|edge} contains the edge that has been reverted.

\item{\bf SET\_STARTNODE}
  
  Is invoked after a node has been marked as being a start-node.
  
  Is delivered to all graph labels and to all labels of the new
  start-node.  {\bf data\verb|->|node} contains the new start-node.

\item{\bf DEL\_STARTNODE}
  
  Is invoked before a node will be marked as no longer being a
  start-node.
  
  Is delivered to all graph labels and to all labels of the
  former start-node.  {\bf data\verb|->|node} contains the former
  start-node.

\item{\bf SET\_ENDNODE}
  
  Is invoked after a node has been marked as being an end-node.
  
  Is delivered to all graph labels and to all labels of the new
  end-node.  {\bf data\verb|->|node} contains the new end-node.

\item{\bf DEL\_ENDNODE}
  
  Is invoked before a node will be marked as no longer being an
  end-node.
  
  Is delivered to all graph labels and to all labels of the
  former end-node.  {\bf data\verb|->|node} contains the former
  end-node.

\item{\bf NAME\_CHANGE}
  
  Is invoked after the name of a graph, node or edge has changed.
  
  Is delivered to either all graph labels, all node labels for
  the node which changed or all edge labels for the edge which
  changes. {\bf data} contains the according node or edge.

\item{\bf ADD\_NODELABEL}
  
  Is invoked after a new node label has been added to the graph.
  
  Is delivered to all graph labels and to all node labels.  {\bf
    data\verb|->|index} contains the index of the new label.

\item{\bf ADD\_EDGELABEL}
  
  Is invoked after a new edge label has been added to the graph.
  
  Is delivered to all graph labels and to all edge labels.  {\bf
    data\verb|->|index} contains the index of the new label.

\item{\bf ADD\_GRAPHLABEL}
  
  Is invoked after a new graph label has been added to the graph.
  
  Is delivered to all graph labels, to all edge labels, and to all
  node labels.  {\bf data\verb|->|index} contains the index of
  the new label.

\item{\bf DEL\_NODELABEL}
  
  Is invoked before a node label will be removed from the graph.
  
  Is delivered to all graph labels and to all node labels,
  including the one that will be removed.  {\bf
    data\verb|->|index} contains the index of the label that will
  be removed.

\item{\bf DEL\_EDGELABEL}
  
  Is invoked before an edge label will be removed from the graph.
  
  Is delivered to all graph labels and to all edge labels,
  including the one that will be removed.  {\bf
    data\verb|->|index} contains the index of the label that will
  be removed.

\item{\bf DEL\_GRAPHLABEL}
  
  Is invoked before an graph label will be removed from the
  graph.
  
  Is delivered to all graph labels, to all node labels, and to all
  edge labels, including the one that will be removed. {\bf
    data\verb|->|index} contains the index of the label that will
  be removed.

\item{\bf SOURCE\_CHANGE}
  
  Is invoked after the source node of an edge changed.
  
  Is delivered to all labels of the edge that changed. {\bf
    data\verb|->|edge} contains the edge that changed.

\item{\bf TARGET\_CHANGE}
  
  Is invoked after the source node of an edge changed.
  
  Is delivered to all labels of the edge that changed. {\bf
    data\verb|->|edge} contains the edge that changed.

\item{\bf DIRECTEDNESS\_CHANGE}
  
  Is invoked after the directedness of the graph changed.
  
  Is delivered to all graph labels, to all node labels, and to all
  edge labels.

\end{itemize}

\end{itemize}

If you write a graph label and want to provide an inspector for the
label you have to assgin the inspector to the {\bf inspector}
member variable. This is usually done in the constructor of the label.
The label must be passed to the constructor of the inspector.

The following methods check if the member variable {\bf inspector !=
NULL} to determine their behaviour or return value:

\begin{itemize}

\item{\bf BOOL FFGraphLabel::hasOwnInspector()}

Returns {\bf YES} to indicate that the the label has its own
interpreter if {\bf inspector != NULL}

\item{\bf void FFGraphLabel::showInspector()}

Displays the inspector on the screen if {\bf inspector !=
NULL}. Usually by calling the method {\bf void FFInspector::show()}.

\item{\bf void FFGraphLabel::updateInspector()}

Updates the display of the inspector if {\bf inspector !=
NULL}. Usually by calling {\bf void FFInspector::inspect()}

\end{itemize}
